% TOP_PROBLEM: {"id": 5500031, "title": "Trapping Light Rays with Segment Mirrors", "category_id": 6, "category": {"id": 6, "name": "geometry", "display_name": "Geometry", "description": "Euclidean and non-Euclidean geometry, geometric structures.", "slug": "geometry", "order_index": 6, "created_at": "2026-07-31T15:26:25.671Z"}, "status": "open", "problem_number": "AMR-054-0031", "set_id": 13}
% FIRST_POSTED: 2026-10-01
% ENTRY_KIND: statement_only
% UnsolvedMath ID 5500031; source catalogue code AMR-054-0031
% !TeX program = lualatex
% Definition and sources only; this file is not an LLM solution attempt.
\documentclass[11pt,a4paper]{article}
\usepackage[margin=25mm]{geometry}
\usepackage{fontspec}
\setmainfont{lmroman10-regular.otf}[BoldFont=lmroman10-bold.otf,
 ItalicFont=lmroman10-italic.otf,BoldItalicFont=lmroman10-bolditalic.otf]
\usepackage{amsmath,amssymb,mathtools}
\usepackage{unicode-math}
\setmathfont{latinmodern-math.otf}
\AtBeginDocument{\let\setminus\smallsetminus}
\usepackage{xurl,hyperref}
\hypersetup{colorlinks=true,urlcolor=blue}
\setcounter{secnumdepth}{0}
\setlength{\parindent}{0pt}
\setlength{\parskip}{5pt}
\setlength{\emergencystretch}{3em}
\begin{document}
\section*{Trapping Light Rays with Segment Mirrors}\label{5500031}
{\small UnsolvedMath ID 5500031 · AMR-054-0031 · Geometry · AMR Open Problem Lists}
\subsection{Definitions and mathematical statement}
Choose finitely many line segments in $\mathbb R^2$ whose closures are
pairwise disjoint. Their relative interiors are two-sided mirrors;
their endpoints are not reflecting. Let $s$ be a point off the mirrors,
and let $H$ be the convex hull of all mirror closures.

A ray from $s$ travels along straight lines. At a transverse hit in
the interior of a mirror it obeys specular reflection: the tangential
velocity component is preserved and the normal component changes sign.
Passing through a missing endpoint gives no reflection. A ray is
trapped when no point of its forward trajectory lies strictly outside
$H$; a ray escaping after a very large but finite number of reflections
is not trapped.

Can the mirrors and source be chosen so that every initial direction
from $s$ gives a trapped ray? The conjectured answer is no: for every
finite arrangement and source there is an escaping direction. It
suffices to consider $s\in H$, since otherwise trapping already fails
at the source.

The quantifier over all directions is essential. Trapping selected
periodic rays, or even a finite collection of nonperiodic rays, does
not give a trap for a point source emitting in every direction.
\subsection{Short English statement}
Can finitely many separated straight mirrors trap every ray emitted
by one point, or must some direction eventually escape?
\subsection{Sources}
\begin{itemize}
\item[S1] ULAM.AI and the UnsolvedMath Contributors, supplied problems.json, record 5500031 (AMR-054-0031), AMR Open Problem Lists. Catalogue identity and source transcription; CC BY 4.0.
\url{https://huggingface.co/datasets/ulamai/UnsolvedMath}.
\item[S2] Open Problems Project - Computational Geometry, Problem 31. Original source indexed by AMR-054-0031.
\url{https://topp.openproblem.net/p31}.
\end{itemize}
\end{document}
